\ssscp{Final consonant}{14-04-03-04}
\ili{As} far as we can see,
\citet{Blust1982} does not justify final *r in **mans(aə)r.
For nearly all our subgroups *R is sufficient to account for the reflexes.
It accounts for the \tsc{\ili{Ambon}-Seram}
and \tsc{Seram-Tanimbar-Bomberai} data (\srf{07-04-01-02}, \srf{10-01}, \srf{11-05}).
It also accounts for \ili{Ambelau} \it{mate} ‘cuscus’ (cf.
*timuR > \it{rimu} ‘\ili{east} monsoon’,
*niuR > \it{niβe} ‘coconut’,
*wahiyəR > \it{βae} ‘water’).
Among \tsc{Aru} languages, \ili{Ujir} \it{meday}
and West \ili{Kola} \it{modah} point to final *j/*R (\srf{06-01-01}).

Among \tsc{Timor-Babar} languages, \ili{Meto},
\ili{Tetun}, Waima{\Q}a, and Ili{\Q}uun all have expected *R > {\0}.\footnote{
		Ili{\Q}uun only has *R > {\0} word finally.}
Dadu{\Q}a has expected final *R > \it{h}.
\ili{Roma}, East \ili{Damer} and \ili{Tugun} have expected *R > \it{r}.
\tsc{Babar} languages (SE. \ili{Babar}, E. \ili{Marsela}, \ili{Emplawas}, Tela)
have *R > \it{r, {\0}} in other words,
so *R > {\0} in these languages is constistent.

Other \tsc{Timor-Babar} languages present several difficulties.
\ili{Galolen} has irregular *R > \it{r} where we expect *R > {\0} (\srf{04-05}),
while \ili{Makuva} has irregular *R > {\0} where we expect *R > \it{r}
(\srf{04-06-01}). 
\ili{Leti} \it{mada}, \ili{Moa} \it{mada},
and \ili{Kisar} \it{masla} usually have *R > \it{r} including in
word-final position, thus requiring apparent irregular *R > {\0} in \ili{Leti}
and \ili{Moa} and *R > \it{l} in \ili{Kisar}.
All these languages have productive processes of final VC{\#}
→ CV{\#} metathesis which is often followed by assimilation and/or
mutation of the original word-final consonant: e.g.
\ili{Leti} \it{βu\tbr{l}a\tbr{n} → βu\tbr{ln}a {\tl} βu\tbr{ll}a}
‘moon’, \it{lɔ\tbr{d}a\tbr{n} → lɔ\tbr{dn}a {\tl} lɔ\tbr{nn}a} ‘rattan’.
(Something similar is attested in \ili{Roma} \it{ma\tbr{d}a\tbr{r} → ma\tbr{tt}a} ‘cuscus’).
Consider also *qatəluR > \it{ternu} ‘egg’,
*kuluR > \it{urnu} ‘bread fruit’,
and *timuR > \it{tipru} ‘\ili{east}’ each of which show unexpected
*l > \it{n} and *m > \it{p} after metathesis (*l = \it{l}
and *m = \it{m} would be expected).
Thus the pathway **mantəR > **mandaR > **madar > **madra > \ili{Leti}
\it{mada} ‘cuscus’ is possible,
though, unfortunately, there are no words of a comparable shape to check.
We note that the consonant cluster \it{dr} does not occur in
any native words in the \ili{Leti} word list in \citet[409--448]{vanEngelenhoven2004}.

\ili{Kisar} \it{masla} appears to show irregular *R > \it{l}.
But the data in \trf{14-10} for ‘rattan’ show an \it{r} > \it{l}
correlation in potential loans such as from \ili{Malay} \it{rotan}.
Consider also \ili{Malay} \it{sə\tbr{r}omboŋ} ‘tube,
cylinder’, \ili{Makasar} \it{sa\tbr{r}omboŋ} ‘pillow case’,
\ili{Tii} (Rote) \it{\tbr{r}ombo} ‘sarong’,
\ili{Kisar} \it{sapo\tbr{l}o} ‘trousers’ (metathesis),
\ili{Luang} \it{po\tbr{l}a} ‘pants’ which also shows an \it{r : l} correspondence.\footnote{
		Data resulting from three-way email exchanges between
		Grimes, Edwards and Antoinette Schapper in September 2020.}
So \it{r > l} seems to be a loan adaptation strategy in \ili{Kisar}
and \ili{Luang}, and has some support even if the \ili{Kisar} form is not
a loan.

For the most part *R accounts for the data we see in the Wallacea
area, while *r leaves many unanswered questions,
being difficult to confirm or dispute,
particularly in word-final position.
Recall also from \srf{14-04-03-01} that final *R explains the
final vowel in the \tsc{\ili{Muna}-Buton} forms.